\section{Security Considerations}


\subsection{Byzantine Fault Tolerance}

Q-Chain maintains safety under standard Byzantine assumptions:

\begin{theorem}[Safety]
If $f < n/3$ validators are Byzantine and the network delay $\Delta < \Delta_{max}$, then no two honest validators finalize conflicting blocks.
\end{theorem}

\begin{proof}
For a block to finalize, it requires:
\begin{enumerate}
\item BLS signatures from $\geq 2n/3$ validators
\item Corona shares from $\geq 2n/3$ validators
\item Confidence $\geq \beta$ in Lux voting
\end{enumerate}

With $f < n/3$ Byzantine nodes, at least $n - f > 2n/3$ honest nodes exist. Two conflicting blocks cannot both obtain $2n/3$ signatures from honest validators.
\end{proof}

\subsection{Liveness}

\begin{theorem}[Liveness]
If $f < n/3$ validators are Byzantine and network delay $\Delta < \Delta_{max}$, then all valid transactions eventually finalize.
\end{theorem}

\begin{proof}
Honest validators always prefer valid transactions. With $>2n/3$ honest validators and bounded network delay, the Lux consensus mechanism guarantees that valid preferences reach confidence threshold $\beta$ within finite rounds.
\end{proof}

\subsection{Slashing Conditions}

\begin{table}[h]
\centering
\begin{tabular}{@{}lll@{}}
\toprule
\textbf{Reason} & \textbf{Evidence} & \textbf{Penalty} \\
\midrule
Double Sign & Two conflicting block sigs & 100\% stake \\
Missing PQ Cert & No Corona signature & 50\% stake \\
Invalid Signature & Malformed signature & 75\% stake \\
Extended Downtime & >99\% missed blocks & 25\% stake \\
\bottomrule
\end{tabular}

\label{tab:slashing}
\end{table}

